\chapter{مروری بر پژوهش‌های پیشین}

\section{مقدمه}	
پیش‌بینی دقیق عملکرد شیر در گاوهای شیری یکی از موضوعات کلیدی در مدیریت بهینه گله‌ها، برنامه‌ریزی تولید، و افزایش سودآوری مزارع لبنی محسوب می‌شود. این فرآیند به دامداران امکان می‌دهد تا با اتخاذ تصمیمات آگاهانه در زمینه تغذیه، سلامت، و انتخاب نژادی دام، بهره‌وری را بهبود بخشند \cite{Dematawewa2007}. در دهه‌های اخیر، با افزایش حجم داده‌های تولیدشده در دام‌پروری و پیشرفت‌های چشمگیر در حوزه تحلیل داده‌ها و یادگیری ماشین، روش‌های پیش‌بینی نیز تحولات قابل‌توجهی را تجربه کرده‌اند \cite{Adriaens2021}. پژوهش حاضر با هدف توسعه مدلی مبتنی بر یادگیری ماشین برای پیش‌بینی اولین رکورد شیردهی در دوره شیردهی بعدی، با استفاده از داده‌های تاریخی و پارامترهای کلیدی از جمله اطلاعات روز تست، تلاش دارد ابزاری نوین برای ارزیابی مدیریت گله، به‌ویژه در دوره انتقال، ارائه دهد. دوره انتقال، به‌عنوان یکی از حساس‌ترین مراحل زندگی گاوهای شیری (شامل سه هفته قبل تا یک ماه پس از زایمان)، نیازمند توجه ویژه‌ای است، زیرا مشکلات متابولیکی و سلامتی در این دوره می‌توانند تأثیرات بلندمدتی بر تولید شیر و سلامت کلی دام داشته باشند \cite{Drackley1999, LeBlanc2010}.
این فصل به بررسی جامع پژوهش‌های پیشین در زمینه پیش‌بینی تولید شیر، کاربرد روش‌های نوین مانند یادگیری ماشین در دام‌پروری، عوامل مؤثر بر عملکرد شیر در اوایل شیردهی، و اهمیت پایش دوره انتقال با استفاده از شاخص‌های خاص مانند شاخص دوره انتقال می‌پردازد. مرور این مطالعات نه‌تنها زمینه علمی پژوهش را فراهم می‌کند، بلکه نقاط قوت و ضعف روش‌های موجود را نیز مشخص می‌سازد تا مسیر توسعه مدل‌های جدید هموار شود. با توجه به تنوع داده‌ها و پیچیدگی‌های بیولوژیکی، این فصل تلاش دارد با ارائه یک دید کلی و عمیق، پایه‌ای محکم برای تحقیقات آتی ایجاد کند. همچنین، با در نظر گرفتن محدودیت‌های موجود در مطالعات پیشین، مانند کمبود داده‌های سلامتی یا عدم تعمیم‌پذیری مدل‌ها، راهکارهایی برای بهبود پیشنهاد خواهد شد.
	
\section{دوره انتقال و اهمیت پایش آن در گاوهای شیری}
دوره انتقال در گاوهای شیری، که به طور معمول بازه شش هفته‌ای در محدوده زایمان را شامل می‌شود، از حساس‌ترین مراحل فیزیولوژیکی در چرخه زندگی تولیدی دام است \cite{Drackley1999}. در این مدت کوتاه، دام تغییرات متابولیکی و هورمونی عمیقی را تجربه می‌کند که شامل کاهش مصرف ماده خشک و همزمان افزایش شدید نیازهای غذایی برای تولید شیر است \cite{Dallago2019}. این عدم تطابق بین مصرف خوراک و نیازهای انرژی، منجر به بروز تعادل منفی انرژی\LTRfootnote{Negative Energy Balance (NEB)} و اختلال در عملکرد سیستم ایمنی می‌شود. در نتیجه، گاوها به شدت مستعد ابتلا به انواع بیماری‌های متابولیکی و عفونی، از جمله کتوز، ورم پستان، متریت، و جابجایی شیردان می‌گردند \cite{LeBlanc2010}. آمارها نشان می‌دهند که ۳۰ تا ۵۰ درصد از گاوها در این دوره دچار حداقل یک اختلال سلامتی می‌شوند که می‌تواند به طور مستقیم بر تولید شیر جاری و عملکرد تولیدمثلی آینده آن‌ها تأثیر منفی بگذارد \cite{Salamone2022}.

روش‌های سنتی برای شناسایی و مدیریت مشکلات دوره انتقال، غالباً بر پایه آزمایش‌های آزمایشگاهی نمونه‌های خون یا شیر (مانند اندازه‌گیری نشانگرهای متابولیکی) استوار هستند \cite{DeKoster2019}. اگرچه این روش‌ها اطلاعات ارزشمندی ارائه می‌دهند، اما اغلب نیازمند مداخله دستی دامدار هستند که می‌تواند به شناسایی دیرهنگام مشکلات منجر شود و کارایی پایش را کاهش دهد \cite{Salamone2022}. این امر، نیاز به توسعه ابزارهای پایش خودکار و داده‌محور را برای شناسایی زودهنگام دام‌های در معرض خطر و اتخاذ تصمیمات پیشگیرانه، بیش از پیش ضروری می‌سازد.
	
\subsection{شاخص دوره انتقال  و محدودیت‌های آن}
یکی از ابزارهای مهمی که برای ارزیابی موفقیت یا شکست دوره انتقال در گاوهای شیری چند شکم‌زا پیشنهاد شده است، شاخص دوره انتقال است \cite{Nordlund2006}. این شاخص بر اساس تفاوت بین تولید شیر مشاهده شده و تولید شیر پیش‌بینی شده در اولین روز تست شیردهی پس از زایمان عمل می‌کند. مدل  شاخص دوره انتقال برای پیش‌بینی تولید شیر، از پارامترهای حاصل از شیردهی جاری و قبلی استفاده می‌کند، از جمله روزهای شیردهی در اولین روز تست، تولید شیر 305	روزه قبلی، طول شیردهی قبلی، نمره لگاریتمی سلول‌های سوماتیک\LTRfootnote{SCC} در آخرین تست شیردهی قبلی، روزهای خشک، و دفعات شیردوشی جاری \cite{Nordlund2006}. با این حال، کاربرد شاخص دوره انتقال عمدتاً به گاوهایی محدود می‌شود که در دومین یا شیردهی‌های بعدی خود قرار دارند، زیرا این شاخص برای عملکرد تلیسه‌ها\LTRfootnote{Heifer} (گاوهای شکم اول) که سابقه شیردهی قبلی ندارند، طراحی نشده است \cite{Dallago2019}.
	
\subsection{چالش پایش دوره انتقال در تلیسه‌ها}	
دوره انتقال برای تلیسه‌ها نیز به همان اندازه گاوهای بالغ چالش‌برانگیز است \cite{Dallago2019}. تلیسه‌ها علاوه بر شروع تولید شیر، همچنان در مرحله رشد بدنی نیز قرار دارند که نیازهای تغذیه‌ای آن‌ها را پیچیده‌تر می‌کند \cite{Dallago2019}. مطالعات نشان داده‌اند که تلیسه‌ها نیز پس از زایمان مستعد اسیدوز شکمبه و کتوز هستند \cite{Dallago2019}. با توجه به محدودیت شاخص دوره انتقال برای تلیسه‌ها، نیاز به توسعه ابزارهایی برای ارزیابی و پیش‌بینی عملکرد شیردهی در این دسته از دام‌ها، به منظور کمک به دامداران در فرآیند تصمیم‌گیری، به شدت احساس می‌شود. این امر به ویژه از آن جهت حائز اهمیت است که یک استراتژی مدیریتی موفق برای گاوهای چند شکم‌زا می‌تواند به طور مشابه برای تلیسه‌ها نیز مؤثر باشد، به شرط آنکه ابزارهای پایش مناسب در دسترس باشند \cite{Dallago2019}.
	
\section{روش‌های پیش‌بینی عملکرد شیر}

\subsection{روش‌های سنتی و آماری}	
پژوهش‌های اولیه در حوزه پیش‌بینی تولید شیر عمدتاً بر پایه مدل‌های آماری کلاسیک استوار بودند. این مدل‌ها، مانند رگرسیون خطی چندگانه\LTRfootnote{Multiple Linear Regression (MLR)} و مدل‌های گاوسین\LTRfootnote{Gaussian Models}، با استفاده از متغیرهایی نظیر مشخصات ظاهری، نژاد، و رکوردهای اولیه شیردهی، تلاش می‌کردند عملکرد آتی دام را تخمین بزنند. به طور سنتی، داده‌های تولید شیر در روزهای تست\LTRfootnote{Test Day(TD)} به منظور پیش‌بینی تولید آتی دام بر اساس منحنی شیردهی مورد استفاده قرار می‌گرفتند. در بسیاری از این موارد، تابع گامای ناقص که توسط وود در سال 1967 پیشنهاد شد، یکی از متداول‌ترین توابع برای توصیف منحنی شیردهی بوده است \cite{Wood1967}. همچنین، مدل‌هایی برای ارزیابی ژنتیکی گاوهای شیری به عنوان جایگزینی برای عملکرد سنتی 305 روزه شیردهی پیشنهاد شده‌اند که قادرند اثرات محیطی را در روز ثبت شیر لحاظ کنند \cite{Ali1987}.
با این حال، این روش‌ها، هرچند پایه‌گذار تحلیل‌های آماری در این حوزه بودند، اما در مدل‌سازی روابط پیچیده و غیرخطی موجود در داده‌های بیولوژیکی دام‌پروری با محدودیت‌هایی مواجه بودند \cite{GuevaraEscobar2022}. به عنوان مثال، \textbf{ گیوارا-اسکوبار و همکاران} در مطالعه خود برای پیش‌بینی تولید شیر 305 روزه از داده‌های روزهای تست، به محدودیت‌هایی در دقت مدل‌های سنتی و وابستگی به داده‌های منظم اشاره کرده‌اند \cite{GuevaraEscobar2022}. همچنین، این مدل‌ها غالباً قادر به شناسایی الگوهای ظریف و پیچیده‌ای که در پس داده‌های بزرگ و متنوع نهفته‌اند، نبودند.
	
\subsection{کاربرد یادگیری ماشین در پیش‌بینی عملکرد شیر}	
با گسترش داده‌های حجیم در دام‌پروری و پیشرفت الگوریتم‌های یادگیری ماشین، رویکردهای نوین جایگزین روش‌های سنتی شدند. الگوریتم‌هایی نظیر درخت تصمیم، شبکه‌های عصبی مصنوعی، ماشین‌های بردار پشتیبانی، و جنگل‌های تصادفی به دلیل توانایی بالای خود در کشف الگوهای پیچیده و مدل‌سازی روابط غیرخطی، به‌طور گسترده‌ای مورد استفاده قرار گرفته‌اند \cite{GuevaraEscobar2022}. استفاده از مدل‌های یادگیری ماشین مزایای قابل توجهی نسبت به مدل‌های رگرسیون سنتی دارد؛ از جمله قابلیت مدیریت داده‌های ناقص، روال‌های اعتبارسنجی متقاطع، انتخاب ویژگی‌ها، بهینه‌سازی هایپرپارامترها، و مدل‌های ترکیبی\LTRfootnote{Ensembel} خودکار که همگی به منظور جلوگیری از بیش‌برازش و بهبود تعمیم‌پذیری مدل طراحی شده‌اند \cite{GuevaraEscobar2022}.
	
\subsection*{پیش‌بینی اولین شیردهی روز تست  با یادگیری ماشین}	
پیش‌بینی اولین رکورد شیردهی\LTRfootnote{First Test Day Milk(FT-MILK)} در شیردهی بعدی، به ویژه در اوایل دوره انتقال، از اهمیت خاصی برخوردار است، چرا که مشکلات سلامتی در این دوره می‌توانند تأثیرات منفی قابل توجهی بر تولید شیر داشته باشند \cite{Dallago2019}.
	\textbf{سالامون و همکاران} به منظور پیش‌بینی اولین رکورد شیردهی در گاوهای شیری، از داده‌های تاریخی روز تست\LTRfootnote{Test Day Records(TDRs)}شیردهی قبلی  و اطلاعات اضافی گاو و گله استفاده کردند \cite{Salamone2022}.
آن‌ها مجموعه‌ای از مدل‌های رگرسیون جنگل تصادفی را توسعه دادند و عملکرد آن‌ها را در برابر چهار مدل پایه\LTRfootnote{Benchmark} ساده‌تر مقایسه کردند. مدل‌های جنگل تصادفی عملکرد پیش‌بینی بهتری داشتند و ریشه میانگین مربع خطا\LTRfootnote{RMSE} آن‌ها بین \num{6,08} تا \num{6,24} کیلوگرم شیر بود، که در محدوده حداقل ریشه میانگین مربع خطا نظری (6 کیلوگرم) برای پیش‌بینی تولید شیر روزانه قرار می‌گیرد \cite{Salamone2022}. این مطالعه نشان داد که استفاده از داده‌های تاریخی تولید شیر امکان پیش‌بینی دقیق تولید در آغاز شیردهی بعدی را فراهم می‌کند و می‌تواند به عنوان بخشی از ابزارهای پایش سلامت مبتنی بر داده مورد استفاده قرار گیرد. مهمترین ویژگی‌های شناسایی شده در این مدل‌ها شامل روزهای شیردهی در اولین روز تست شیردهی بعدی\LTRfootnote{DIM TD1X+1} و تولید شیر 305 روزه تجمعی\LTRfootnote{M305} شیردهی قبلی بود. همچنین، تولید شیر در روزهای تست چهارم و پنجم شیردهی قبلی نیز اهمیت بالایی در پیش‌بینی داشت \cite{Salamone2022}.
	
\textbf{دالاگو و همکاران} با تمرکز بر پیش‌بینی اولین رکورد شیردهی در تلیسه‌های شیری، به این نتیجه رسیدند که مدل‌های یادگیری ماشین پتانسیل بالایی برای پر کردن خلاء موجود در ابزارهای پایش دوره انتقال برای این دسته از دام‌ها دارند \cite{Dallago2019}.
آن‌ها سه نوع مدل رگرسیون خطی چندمتغیره، جنگل تصادفی و شبکه عصبی مصنوعی را ارزیابی کردند. نتایج نشان داد که مدل شبکه عصبی مصنوعی عملکرد بهتری در پیش‌بینی اولین رکورد شیردهی نسبت به رگرسیون خطی چند‌متغیره و جنگل تصادفی دارد (میانگین خطای مطلق کمتر از 4 کیلوگرم). 
این مطالعه همچنین بر اهمیت متغیرهایی مانند وزن در زمان زایمان، روزهای شیردهی در روز تست، و بتا-هیدروکسی‌بوتیرات\LTRfootnote{Beta-hydroxybutyrate} شیر (به عنوان نشانگر کتوز) در پیش‌بینی اولین رکورد شیر تأکید کرد. در حالی که مدل جنگل تصادفی دچار بیش‌برازش شد، مدل شبکه عصبی مصنوعی پایداری و تعمیم‌پذیری بهتری نشان داد که آن را به گزینه‌ای مناسب‌تر برای مدل‌سازی معیارهای طرح بهبود گله‌های شیری\LTRfootnote{Dairy Herd Improvement(DHI)} تبدیل می‌کند \cite{Dallago2019}. این تحقیق اولین قدم مهمی در توسعه ابزاری مشابه شاخص دوره انتقال برای تلیسه‌ها محسوب می‌شود.
	
\subsection*{پیش‌بینی تولید شیر 305 روزه با داده‌های اولیه شیردهی}
در حالی که پژوهش حاضر بر پیش‌بینی اولین رکورد شیردهی متمرکز است، پژوهش‌هایی که به پیش‌بینی کل دوره شیردهی (مانند 305 روزه) با استفاده از داده‌های اولیه شیردهی می‌پردازند، نیز ارتباط نزدیکی با هدف این تحقیق دارند.
\textbf{گیوارا-اسکوبار و همکاران} به دنبال دستیابی به یک مدل یادگیری ماشین برای پیش‌بینی تولید شیر 305	روزه\LTRfootnote{MY305} از همان دوره شیردهی، با استفاده از داده‌های روزهای تست بودند \cite{GuevaraEscobar2022}. 
آن‌ها پایگاه داده‌ای از 11,892 رکورد تولید شیر روزانه از گاوهای هلشتاین-فریزین\LTRfootnote{Holstein-Friesian} را مورد استفاده قرار دادند. مدل‌های یادگیری ماشین مختلفی از جمله یادگیری عمیق، ماشین تقویت گرادیان\LTRfootnote{Gradient Boost Machine(GBM)}، افزایش گرادیان افراطی و مدل خطی تعمیم‌یافته مورد بررسی قرار گرفتند. نتایج آن‌ها نشان داد که مدل‌های ترکیبی بهترین برازش را در پیش‌بینی شیر تجمعی 305 روز ارائه می‌دهند (ضریب تعیین \num{0,79} و ریشه میانگین مربع خطا 1256).

این مطالعه تأکید کرد که تولید شیر در فاز اولیه شیردهی (روزهای 5,10,20,30) مهمترین متغیرها در مدل‌های پیش‌بینی شیرتجمعی 305 روز بودند. این رویکرد نشان می‌دهد که حتی با داده‌های محدود از روزهای اولیه شیردهی، می‌توان پیش‌بینی‌های دقیقی از عملکرد کلی شیردهی به دست آورد که می‌تواند برای اتخاذ تصمیمات مدیریتی مهم در تغذیه، سلامت، و بهبود ژنتیکی به کار رود \cite{GuevaraEscobar2022}.
	
\section{عوامل تأثیرگذار بر عملکرد شیر در اوایل شیردهی و دوره انتقال}

مطالعات گوناگون نشان داده‌اند که عوامل متعددی بر عملکرد شیر در اوایل دوره شیردهی و به طور کلی بر موفقیت دوره انتقال تأثیرگذار هستند. شناخت این عوامل برای ساخت مدل‌های پیش‌بینی دقیق بسیار حیاتی است:
	\begin{enumerate}
	\item \textbf{سابقه تولید شیر پیشین:} داده‌های تاریخی مانند تولید 305 روزه شیردهی قبلی\LTRfootnote{M305X}، تولید شیر در روزهای تست خاص در شیردهی قبلی (مانند روزهای تست چهارم، پنجم و ششم) و حداکثر تولید شیر در شیردهی قبلی، از مهم‌ترین ورودی‌ها برای پیش‌بینی عملکرد آتی دام هستند \cite{Salamone2022}. این سوابق به مدل‌ها امکان می‌دهند تا ظرفیت تولیدی ژنتیکی و سابقه عملکردی هر دام را در نظر بگیرند.
	
	\item \textbf{روزهای شیردهی در زمان تست\LTRfootnote{Days In Milk(DIM)}:} این متغیر نقش مستقیمی در تعیین مقدار شیر ثبت‌شده در اولین تست دارد و به دلیل شیب صعودی منحنی شیردهی در آغاز شیردهی، تأثیر قابل توجهی بر پیش‌بینی اولین رکورد شیردهی دارد \cite{Salamone2022, Dallago2019}. همچنین، تولید شیر در روزهای اولیه شیردهی (مانند روزهای 5,10,20,30) به عنوان مهمترین متغیرها در پیش‌بینی شیر تجمعی 305 روزه شناخته شده‌اند \cite{GuevaraEscobar2022}.
	
	\item \textbf{دوره شیردهی:} الگوی منحنی شیردهی و پتانسیل تولید شیر گاو بسته به شماره زایش دام تغییر می‌کند. گاوهای جوان‌تر (تلیسه‌ها) ممکن است الگوی رشد متفاوتی نسبت به گاوهای بالغ داشته باشند و این عامل باید در مدل‌سازی لحاظ شود \cite{GuevaraEscobar2022, Dallago2019}.
	
	\item \textbf{طول دوره خشک و فاصله گوساله‌زایی:} مدت زمان دوره خشک بین دو شیردهی و فاصله گوساله‌زایی نیز بر تولید شیر بعدی تأثیرگذار است. یک دوره خشک بهینه برای بازسازی بافت پستان و آماده‌سازی برای شیردهی بعدی ضروری است \cite{Salamone2022}.
	
	\item \textbf{سلامت پستان و شمارش سلول‌های سوماتیک\LTRfootnote{SCC}:} شمارش بالای سلول‌های سوماتیک معمولاً با التهاب پستان (ورم پستان) و کاهش عملکرد شیر همراه است. این شاخص می‌تواند به عنوان یک نشانگر مهم برای سلامت پستان و در نتیجه تأثیرگذاری بر تولید شیر در نظر گرفته شود \cite{GuevaraEscobar2022}.
	
	\item \textbf{شاخص‌های متابولیکی و بهداشتی:} نشانگرهایی مانند بتا-هیدروکسی‌بوتیرات شیر می‌توانند به عنوان شاخصی برای کتوزیس (یک اختلال متابولیکی شایع در دوره انتقال) عمل کنند و تأثیر منفی قابل توجهی بر تولید شیر اولیه دارند \cite{Dallago2019}. سایر عوامل مانند دوقلوزایی، سهولت زایمان، و اندازه گوساله نیز می‌توانند بر سلامت دام و در نتیجه اولین رکورد شیر تأثیر بگذارند \cite{Dallago2019}.
	
	\item \textbf{وزن دام در زمان زایمان:} وزن مناسب در زمان زایمان با عملکرد تولیدی بهتر در تلیسه‌ها مرتبط است و می‌تواند به عنوان یک عامل مثبت در پیش‌بینی اولین رکورد شیر در نظر گرفته شود \cite{Dallago2019}.
	
	\item \textbf{فاکتورهای محیطی و مدیریتی در سطح گله:} عواملی نظیر ماه زایش یا فصل شیردهی و میانگین‌های عملکردی گله (مانند میانگین تولید شیر تجمعی در روزهای مختلف، میانگین سن اولین زایمان، میانگین فاصله گوساله‌زایی، و میانگین روزهای خشک) نیز می‌توانند تأثیر قابل توجهی بر عملکرد فردی دام داشته باشند. این عوامل نشان‌دهنده شرایط مدیریتی و محیطی حاکم بر گله هستند و می‌توانند تفاوت در سطوح تولید را توجیه کنند \cite{Salamone2022, GuevaraEscobar2022}.
\end{enumerate}	
شناسایی و گنجاندن این متغیرهای کلیدی در مدل‌های پیش‌بینی، دقت و قدرت تفسیر آن‌ها را به طور قابل توجهی افزایش می‌دهد و بینش‌های ارزشمندی را برای بهبود شیوه‌های مدیریتی ارائه می‌کند.
	
\section{چالش‌ها، محدودیت‌ها و پژوهش‌های آتی}
با وجود پیشرفت‌های قابل توجه در زمینه پیش‌بینی عملکرد شیر با استفاده از یادگیری ماشین، همچنان چالش‌ها و محدودیت‌هایی وجود دارد که مسیر را برای پژوهش‌های آتی هموار می‌سازد:
	
\begin{itemize}
	\item \textbf{کیفیت و تکمیل بودن داده‌ها:} یکی از بزرگترین محدودیت‌ها، عدم دسترسی به پایگاه داده‌های جامع و با کیفیت از اطلاعات سلامتی دام است. نبود اطلاعات دقیق بیماری‌ها، انتخاب گاوهای "سالم" و "بدون اختلال" در دوره انتقال را دشوار می‌سازد و می‌تواند بر دقت مدل‌های پیش‌بینی تولید شیر "مورد انتظار" تأثیر بگذارد \cite{Salamone2022}. علاوه بر این، پراکندگی داده‌های روز تست شیردهی و نبود ثبت منظم تمام اطلاعات (مانند درصد چربی، پروتئین، یا شمارش سلول‌های سوماتیک در تمامی تست‌ها) می‌تواند کارایی مدل‌ها را کاهش دهد \cite{GuevaraEscobar2022}.
	
	\item \textbf{تعمیم‌پذیری مدل‌ها:} بسیاری از مطالعات، از جمله \textbf{دالاگو و همکاران}، تنها بر یک نژاد خاص (مانند هلشتاین) متمرکز بوده‌اند \cite{Dallago2019}. این امر نیاز به توسعه و ارزیابی مدل‌ها برای نژادهای دیگر را ضروری می‌سازد تا قابلیت تعمیم‌پذیری آن‌ها در سناریوهای مختلف دامداری تضمین شود. همچنین، تغییرات قابل توجه بین گله‌ها (مانند تفاوت در ژنتیک گله، شیوه‌های مدیریتی، و محیط) می‌تواند بر عملکرد مدل‌ها تأثیر بگذارد و گنجاندن "فاکتور گله" به عنوان یک متغیر در مدل‌ها، می‌تواند دقت پیش‌بینی را بهبود بخشد \cite{Dallago2019}.
	
	\item \textbf{بیش‌برازش در مدل‌ها:} در برخی موارد، مانند مدل جنگل تصادفی در مطالعه \textbf{دالاگو و همکاران}، پدیده بیش‌برازش مشاهده شده است که می‌تواند توانایی تعمیم‌پذیری مدل را به داده‌های جدید کاهش دهد \cite{Dallago2019}. اگرچه الگوریتم‌های یادگیری ماشین مدرن دارای مکانیزم‌هایی برای مقابله با بیش‌برازش هستند، اما انتخاب صحیح هایپرپارامترها و اعتبارسنجی دقیق مدل برای جلوگیری از این پدیده حیاتی است.
	
	\item \textbf{محدودیت‌های روزهای تست شیردهی:} اگرچه داده‌های روز تست شیردهی به طور گسترده در دسترس هستند و استانداردسازی شده‌اند، اما به دلیل فاصله زمانی بین ثبت رکوردها (معمولاً هر چهار تا هشت هفته)، ممکن است رکوردهای مربوط به اوایل شیردهی بسیار دیرتر از زمان بروز مشکلات سلامت ثبت شوند \cite{Salamone2022}. این محدودیت می‌تواند کاربرد این داده‌ها را در پایش زمان واقعی دوره انتقال با چالش مواجه کند. با این حال، اگر قدرت پیش‌بینی مشکلات دوره انتقال با استفاده از این داده‌ها اثبات شود، می‌توان به فکر تغییر روش ثبت تولید در مراحل اولیه شیردهی بود (مثلاً افزایش دفعات ثبت در هفته‌های اول) \cite{Salamone2022}.
	
	\item \textbf{اعتبار سنجی کاربردی مدل‌ها:} مدل‌های پیش‌بینی‌کننده توسعه یافته، اگرچه عملکرد آماری خوبی از خود نشان داده‌اند، اما نیاز به اعتبارسنجی گسترده در محیط‌های واقعی دامداری دارند تا میزان کارایی آن‌ها به عنوان ابزارهای عملی پایش و تصمیم‌گیری مورد ارزیابی قرار گیرد \cite{Salamone2022, Dallago2019}. این شامل بررسی تأثیر دام‌هایی است که به دلیل بیماری‌های مرتبط با دوره انتقال، پیش از اولین روز تست حذف می‌شوند، بر قابلیت‌های پایش مدل‌ها \cite{Salamone2022}.
\end{itemize}

پژوهش‌های آتی می‌توانند بر جمع‌آوری و یکپارچه‌سازی پایگاه داده‌های غنی‌تر شامل اطلاعات سلامت دقیق‌تر، ارزیابی مدل‌ها بر روی نژادهای مختلف و در محیط‌های گله‌ای متنوع، و توسعه مدل‌هایی با قابلیت تفسیرپذیری بالاتر  تمرکز کنند تا کاربرد عملی این ابزارها در صنعت دامپروری به حداکثر رسد \cite{GuevaraEscobar2022}.
	
\section{نتیجه‌گیری}	
مرور پژوهش‌های پیشین نشان می‌دهد که پیش‌بینی دقیق عملکرد شیر در گاوهای شیری، به‌ویژه در اوایل دوره شیردهی و دوره حساس انتقال، از اهمیت حیاتی برای مدیریت بهینه دامداری برخوردار است. در حالی که روش‌های آماری سنتی محدودیت‌هایی در مدل‌سازی روابط پیچیده دارند، پیشرفت‌های اخیر در یادگیری ماشین راهکارهای قدرتمندی را ارائه کرده‌اند. مطالعات اخیر، از جمله تحقیقات \textbf{سالامون و همکاران}، \textbf{دالاگو و همکاران}، و \textbf{گیوارا-اسکوبار و همکاران}، به وضوح پتانسیل بالای الگوریتم‌هایی نظیر جنگل تصادفی، شبکه‌های عصبی مصنوعی، و مدل‌های ترکیبی را در پیش‌بینی دقیق تولید شیر نشان داده‌اند. این مدل‌ها قادرند با بهره‌گیری از داده‌های تاریخی تولید شیر (مانند رکورد 305	روزه شیردهی قبلی و رکوردهای روز تست اولیه) و سایر اطلاعات کلیدی (نظیر دوره شیردهی، وزن در زمان زایمان، و شاخص‌های سلامتی مانند بتا-هیدروکسی‌بوتیرات)، عملکرد شیردهی آینده دام را با دقت مناسبی پیش‌بینی کنند.
این پژوهش‌ها تأکید می‌کنند که پیش‌بینی اولین رکورد شیردهی در دوره بعدی، می‌تواند به عنوان یک ابزار پایش مؤثر برای شناسایی زودهنگام دام‌های در معرض خطر در دوره انتقال عمل کند و به دامداران امکان دهد تا مداخلات مدیریتی و بهداشتی لازم را به موقع انجام دهند. با این وجود، چالش‌هایی نظیر کیفیت و کامل بودن داده‌های سلامت، تعمیم‌پذیری مدل‌ها به نژادها و گله‌های مختلف، و نیاز به اعتبارسنجی عملی مدل‌های توسعه‌یافته همچنان وجود دارد.
پژوهش حاضر، با تمرکز بر پیش‌بینی اولین رکورد شیر با استفاده از یادگیری ماشین و اطلاعات دوره‌های پیشین و همچنین داده‌های روز تست، در راستای پر کردن این شکاف‌های پژوهشی و توسعه ابزاری کاربردی برای دامداران کشور گام برمی‌دارد. هدف نهایی این تحقیق، ارائه یک مدل قوی و قابل اعتماد است که بتواند به بهبود سلامت گله، افزایش بهره‌وری، و ارتقاء پایداری صنعت دامپروری کشور کمک کند.
	